\documentclass[../../main.tex]{subfiles}% 注意这里的文件路径不能用 ./main.tex ,否则用latexmk编译子文件会报错
\graphicspath{{\subfix{./image/}}} % 指定图片目录,后续可以直接使用图片文件名
% 注意这里的文件路径不能用 ../../image/ ,否则用latexmk编译子文件会报错

% 例如:
% \begin{figure}[H]
% \centering
% \includegraphics[scale=0.3]{图.png}
% \caption{}
% \label{figure:图}
% \end{figure}
% 注意:上述\label{}一定要放在\caption{}之后,否则引用图片序号会只会显示??.

\begin{document}

\section{外形式}
\begin{definition}
设$V$是$n$维线性空间，$V^*$是其对偶空间，**$r$次外形式**是$V$上的反对称$r$重线性函数，全体$r$次外形式构成的线性空间记作$\Lambda^rV^*$。
\end{definition}

\begin{definition}
外形式的**外积**运算$\wedge:\Lambda^rV^*\times\Lambda^sV^*\to\Lambda^{r+s}V^*$满足：
\begin{enumerate}[(1)]
\item 双线性性：$(\lambda\omega_1+\mu\omega_2)\wedge\eta=\lambda\omega_1\wedge\eta+\mu\omega_2\wedge\eta$；
\item 反对称性：$\omega\wedge\eta=(-1)^{rs}\eta\wedge\omega$；
\item 结合律：$(\omega\wedge\eta)\wedge\theta=\omega\wedge(\eta\wedge\theta)$。
\end{enumerate}
\end{definition}

\begin{proposition}
设$\{\omega^1,\omega^2\}$是二维线性空间对偶空间的基，则：
\begin{enumerate}[(1)]
\item $\Lambda^0V^*=\mathbb{R}$；
\item $\Lambda^1V^*=\mathrm{span}\{\omega^1,\omega^2\}$；
\item $\Lambda^2V^*=\mathrm{span}\{\omega^1\wedge\omega^2\}$；
\item $r\geqslant3$时，$\Lambda^rV^*=\{0\}$。
\end{enumerate}
\end{proposition}

\section{外微分}
\begin{definition}
光滑曲面$M$上的**外微分**$\mathrm{d}:\Lambda^r(M)\to\Lambda^{r+1}(M)$是线性映射，满足：
\begin{enumerate}[(1)]
\item 对$f\in\Lambda^0(M)$，$\mathrm{d}f$是$f$的普通微分；
\item $\mathrm{d}(\omega\wedge\eta)=\mathrm{d}\omega\wedge\eta+(-1)^r\omega\wedge\mathrm{d}\eta$；
\item $\mathrm{d}(\mathrm{d}\omega)=0$。
\end{enumerate}
\end{definition}

\begin{example}
设$\omega=P\mathrm{d}u+Q\mathrm{d}v\in\Lambda^1(M)$，则其外微分为
\[\mathrm{d}\omega=\left(\dfrac{\partial Q}{\partial u}-\dfrac{\partial P}{\partial v}\right)\mathrm{d}u\wedge\mathrm{d}v.\]
\end{example}

\section{活动标架}
\begin{definition}
沿黎曼曲面$M$定义的单位正交切向量场$\{e_1,e_2\}$称为**活动标架**，其对偶标架$\{\omega^1,\omega^2\}$满足$\omega^i(e_j)=\delta^i_j$。
\end{definition}

\begin{definition}
活动标架的**联络形式**$\omega^1_2$满足
\[\mathrm{d}\omega^1=-\omega^1_2\wedge\omega^2,\quad \mathrm{d}\omega^2=-\omega^2_1\wedge\omega^1,\quad \omega^2_1=-\omega^1_2.\]
\end{definition}

\section{曲面论的标架形式}
\begin{theorem}
在活动标架下，曲面的高斯曲率满足
\[\mathrm{d}\omega^1_2=-K\omega^1\wedge\omega^2,\]
该式称为**高斯方程**的外微分形式。
\end{theorem}

\begin{proof}
由联络形式的定义和外微分的性质，直接计算可得高斯曲率的外微分表达式。

\end{proof}

\begin{proposition}
曲面的结构方程为
\begin{align*}
&\mathrm{d}\omega^1=-\omega^1_2\wedge\omega^2,\\
&\mathrm{d}\omega^2=-\omega^2_1\wedge\omega^1,\\
&\mathrm{d}\omega^1_2=-K\omega^1\wedge\omega^2.
\end{align*}
\end{proposition}

\section{高斯-博内定理的外微分证明}
\begin{proof}
设$M$是紧致可定向黎曼曲面，取单位正交活动标架$\{e_1,e_2\}$，对偶标架$\{\omega^1,\omega^2\}$，联络形式$\omega^1_2$。
由Stokes公式
\[\iint\limits_M\mathrm{d}\omega^1_2=\int\limits_{\partial M}\omega^1_2,\]
结合$\mathrm{d}\omega^1_2=-K\omega^1\wedge\omega^2$，对曲面进行三角剖分，计算边界积分可得
\[\iint\limits_M K\mathrm{d}\sigma=2\pi\chi(M).\]

\end{proof}

\section{外微分法的应用举例}
\begin{example}
用外微分法证明：高斯曲率$K\equiv0$的曲面局部上与平面等距。
\end{example}

\begin{solution}
由$K\equiv0$得$\mathrm{d}\omega^1_2=0$，根据Poincaré引理，存在函数$\theta$使得$\omega^1_2=\mathrm{d}\theta$。
令
\begin{align*}
&\widetilde{\omega}^1=\cos\theta\omega^1+\sin\theta\omega^2,\\
&\widetilde{\omega}^2=-\sin\theta\omega^1+\cos\theta\omega^2,
\end{align*}
则$\mathrm{d}\widetilde{\omega}^1=\mathrm{d}\widetilde{\omega}^2=0$，故存在局部坐标$u,v$使得$\widetilde{\omega}^1=\mathrm{d}u$，$\widetilde{\omega}^2=\mathrm{d}v$，黎曼度量化为$(\mathrm{d}u)^2+(\mathrm{d}v)^2$，与平面等距。

\end{solution}

\begin{remark}
外微分法简化了曲面论的计算，避免了复杂的协变导数运算，是现代微分几何的核心工具。
\end{remark}




\end{document}